\documentclass[a4paper,11pt]{article}
\usepackage[utf8]{inputenc}
\usepackage[T1]{fontenc}
\usepackage{amsmath, amssymb}
\usepackage{geometry}
\usepackage{parskip}
\usepackage{fancyhdr}
\usepackage{listings}
\usepackage{xcolor}
\usepackage{booktabs}

\geometry{margin=2.5cm}

\pagestyle{fancy}
\fancyhf{}
\rhead{Competitive Programming}
\lhead{Two Sum}
\cfoot{\thepage}

\lstset{
    basicstyle=\ttfamily\small,
    backgroundcolor=\color{gray!10},
    frame=single,
    breaklines=true,
}

\begin{document}

\begin{center}
    {\LARGE \textbf{Two Sum}} \\[0.4em]
    {\large Time Limit: 2.0s \quad Memory Limit: 256 MB}
\end{center}

\hrule
\vspace{1em}

% ─── Problem Statement ────────────────────────────────────────────────────────
\section*{Problem Statement}

You are given an array of integers \texttt{nums} and an integer \texttt{target}.
Return the \textbf{0-based indices} of the two numbers such that they add up to
\texttt{target}.

You may assume that each input has \textbf{exactly one solution}, and you may
\textbf{not} use the same element twice. You may return the answer in any order.

% ─── Input Format ─────────────────────────────────────────────────────────────
\section*{Input Format}

\begin{itemize}
    \item The first line contains two integers: $n$ (the size of the array) and
          $\text{target}$.
    \item The second line contains $n$ space-separated integers $\text{nums}[i]$.
\end{itemize}

% ─── Output Format ────────────────────────────────────────────────────────────
\section*{Output Format}

Print two space-separated integers representing the 0-based indices of the two
numbers that add up to \texttt{target} (in any order).

% ─── Constraints ──────────────────────────────────────────────────────────────
\section*{Constraints}

\[
    2 \le n \le 1000
\]
\[
    -10^9 \le \text{nums}[i] \le 10^9
\]
\[
    -10^9 \le \text{target} \le 10^9
\]

\textbf{Guarantee:} Exactly one valid pair of indices exists.

% ─── Sample ───────────────────────────────────────────────────────────────────
\section*{Sample}

\subsection*{Input}
\begin{lstlisting}
4 9
2 7 11 15
\end{lstlisting}

\subsection*{Output}
\begin{lstlisting}
0 1
\end{lstlisting}

\subsection*{Explanation}
$\text{nums}[0] + \text{nums}[1] = 2 + 7 = 9 = \text{target}$, so we output
\texttt{0 1}.

% ─── Notes ────────────────────────────────────────────────────────────────────
\section*{Notes}

\begin{itemize}
    \item The array contains integers that may be negative or zero.
    \item The answer indices must be \emph{distinct} ($i \neq j$).
    \item There is guaranteed to be exactly one solution; no need to handle the
          case of no solution or multiple solutions.
\end{itemize}

\end{document}
